\section*{Resources and data availability}
The probe implementation, the injection and analysis code, and the per-asset probe statistics can be released with this paper.
3D-DefectBench meshes are not redistributed: the dataset is CC BY-NC 4.0 and was used for non-commercial research only \citep{zhao2026defectbench}.
MATE-3D and Google Scanned Objects are CC BY 4.0 \citep{zhang2025mate3d,downs2022gso}.
Hi3DBench is distributed under the MIT licence on its dataset card \citep{zhang2025hi3deval}.
DB-3DME, as released, provided GIF renders rather than meshes \citep{jia2026db3dme}.
The author works on modelfy.art (\url{https://modelfy.art}), an online image- and text-to-3D service; no modelfy.art models, assets or data were used in this study.

\section*{Declaration of competing interest}
The author is affiliated with modelfy.art, a commercial 3D-generation service, which could benefit from improved evaluation of generated 3D assets.
The study uses only public datasets and open-source tools; no modelfy.art system was evaluated.

\section*{Acknowledgements}
We thank the authors of 3D-DefectBench, MATE-3D, Hi3DBench, DB-3DME and Google Scanned Objects for releasing data, labels and, in the case of 3D-DefectBench, VLM predictions.
Probe code uses Trimesh, PyMeshLab and Embree \citep{trimesh,cignoni2008meshlab,wald2014embree}.
